%------------------------------------------------------------------------------%
% Luis A. D'Afonseca
%
%------------------------------------------------------------------------------%
\documentclass[a4paper,12pt,fleqn]{article}

\usepackage{style_avaliacao}

\ShowAnswers

%------------------------------------------------------------------------------%
\begin{document}

\makeHeader{Integrais e Séries}{Exame -- Turma 6}{25/02/2025}

O exercício correspondente a prova que vai ser substituída vale 40 pontos.

\vspace{\baselineskip}
Tabela de integrais
\begin{align*}
  \int \frac{dx}{a^2+x^2} & = \frac{1}{a}\arctan\left(\frac{x}{a}\right) + c &
  \int \csc^2(x) \, dx & = -\cot(x) + c
\end{align*}

% 01
%------------------------------------------------------------------------------%
%\vspace{\baselineskip}
\exercise{20}
Encontre o volume do sólido obtido pela rotação da região contida entre
os gráficos das funções \(f(x)=\abs{x-1}\) e \(g(x)=\frac{2-x}{2}\)
em torno da reta \(x=-1\)

\begin{answer}
As funções vão se interceptar em $x=0$ e em $b\in[1, 2]$
\begin{align*}
  f(b)            & = g(b)                                             \\
  \abs{b-1}       & = 1 - \frac{b}{2}                                  \\
  b-1             & = 1 - \frac{b}{2} \hspace{20mm} \text{pois } b > 1 \\
  \frac{3}{2} b   & = 2                                                \\
  b = \frac{4}{3} &
\end{align*}

Vamos usar cascas cilíndricas
\[
  V = \int_a^b 2\pi r(x) h(x) dx
\]
onde
\begin{align*}
  r(x) & = x + 1       \\
  h(x) & = g(x) - f(x) \\
  a    & = 0           \\
  b    & = \frac{4}{3}
\end{align*}

Para encontrar $h(x)$ precisamos dividir região em duas partes, para $x\in[0, 1]$
\[
  h(x)
  = g(x) - f(x)
  = 1 -\frac{x}{2} - \abs{x-1}
  = 1 -\frac{x}{2} + x - 1
  = \frac{x}{2}
\]
para \(x\in\left[1, \sfrac{4}{3}\right]\)
\[
  h(x)
  = g(x) - f(x)
  = 1 -\frac{x}{2} - \abs{x-1}
  = 1 -\frac{x}{2} - x + 1
  = 2 -\frac{3x}{2}
\]
Calculando o volume para $x\in[0, 1]$
\begin{align*}
  V_1
  & = \int_0^1 2\pi r(x) h(x) \, dx \\
  & = \int_0^1 2\pi (x+1) \, \frac{x}{2}  \, dx \\
  & = \pi \int_0^1 x^2 + x\,  dx \\
  & = \pi \left(\frac{x^3}{3} + \frac{x^2}{2}\right)\evalat{0}{1}  \\
  & = \pi \left(\frac{1}{3} + \frac{1}{2}\right)  \\
  & = \frac{5}{6}\pi
\end{align*}
para \(x\in\left[1, \sfrac{4}{3}\right]\)
\begin{align*}
  V_2
  & = \int_1^{\sfrac{4}{3}} 2\pi r(x) h(x) \, dx \\
  & = \int_1^{\sfrac{4}{3}} 2\pi (x+1) \, \left(2 -\frac{3x}{2}\right)  \, dx \\
  & = 2\pi \int_1^{\sfrac{4}{3}} -\frac{3}{2}x^2 + \frac{x}{2} + 2 \, dx \\
  & = 2\pi \left(-\frac{x^3}{2} + \frac{x^2}{4} + 2x \right)\evalat{1}{\sfrac{4}{3}}  \\
  & = \pi \left(-x^3 + \frac{x^2}{2} + 4x \right)\evalat{1}{\sfrac{4}{3}}  \\
  & = \pi
  \left[
    \left(-\frac{4^3}{3^3} + \frac{4^2}{3^2\times2} + \frac{4^2}{3}\right)-
    \left(-1 + \frac{1}{2} + 4\right)
   \right] \\
   & = \pi
   \left(-\frac{4^3}{3^3} + \frac{4^2}{3^2\times2} + \frac{4^2}{3} -3 -\frac{1}{2}\right) \\
   & = \frac{19}{54}\pi
\end{align*}
Volume do sólido
\[
  V
  = V_1 + V_2
  = \frac{5}{6}\pi + \frac{19}{54}\pi
  = \frac{32}{27}\pi
  \approx 3,7234
\]
\end{answer}

% 02
%------------------------------------------------------------------------------%
%\vspace{\baselineskip}
\exercise{20}
Calcule a integral
\(
  \int \frac{x^2}{\sqrt{4-x^2}}\,dx
\)

\begin{answer}
  \begin{align*}
    F(x)
    & = \int \frac{x^2}{\sqrt{4-x^2}}\,dx
  \end{align*}
  \[
    x = 2\sin(\theta) \qquad dx = 2\cos(\theta) d\theta
  \]
  \[
    \sqrt{4-x^2}
    = \sqrt{4-2^2\sin^2(\theta)}
    = \sqrt{4\left(1-\sin^2(\theta)\right)}
    = \sqrt{4\cos^2(\theta)}
    = 2\cos(\theta)
  \]
  \begin{align*}
    F(x)
    & = \int \frac{x^2}{\sqrt{4-x^2}}\,dx \\
    & = \int \frac{4\sin^2(\theta)}{2\cos(\theta)}\,2\cos(\theta)\,d\theta \\
    & = 4 \int \sin^2(\theta)\,d\theta \\
    & = 4 \int \frac{1}{2}\left(1-\cos(2\theta)\right)\,d\theta \\
    & = 2 \int 1-\cos(2\theta)\,d\theta \\
    & = 2 \left(\theta - \frac{\sin(2\theta)}{2}\right) + c \\
    & = 2\theta - \sin(2\theta) + c \\
    & = 2\theta - 2\cos(\theta)\sin(\theta) + c \\
    & = 2\arcsin\left(\frac{x}{2}\right) - 2\cos(\theta)\frac{x}{2} + c \\
    & = 2\arcsin\left(\frac{x}{2}\right) - \frac{\sqrt{4-x^2}}{2}x + c
  \end{align*}

\end{answer}

% 03
%------------------------------------------------------------------------------%
%\vspace{\baselineskip}
\exercise{20}
Analise a convergência da série \(\;\sum_{n=2}^\infty \frac{1}{n\left(1+\ln^2(n)\right)}\)

\begin{answer}
Usando o teste da integral com a função
\[
  f(x) = \frac{1}{x\left(1+\ln^2(x)\right)}
\]
Verificando as condições do teste

A função $f$ só não é contínua em $x=0$, portanto é contínua em $[2, \infty)$.

A funçao $f$ é positiva para todo $x>0$.

Derivando $f$ temos
\begin{align*}
  f'(x)
  & = \frac{d}{dx}\left[x\left(1+\ln^2(x)\right)\right]^{-1} \\
  & = -\left[x\left(1+\ln^2(x)\right)\right]^{-2}
       \frac{d}{dx}\left[x\left(1+\ln^2(x)\right)\right] \\
  & = \frac{-1}{x^2\left(1+\ln^2(x)\right)^2}
      \left[\left(1+\ln^2(x)\right) + x\left(1+\ln^2(x)\right)'\right] \\
  & = \frac{-1}{x^2\left(1+\ln^2(x)\right)^2}
      \left[1+\ln^2(x) + x\left(\ln^2(x)\right)'\right] \\
  & = \frac{-1}{x^2\left(1+\ln^2(x)\right)^2}
      \left[1+\ln^2(x) + x2\ln(x)\left(\ln(x)\right)'\right] \\
  & = \frac{-1}{x^2\left(1+\ln^2(x)\right)^2}
      \left[1+\ln^2(x) + x2\ln(x)\frac{1}{x}\right] \\
  & = -\frac{1+\ln^2(x) + 2\ln(x)}{x^2\left(1+\ln^2(x)\right)^2}
\end{align*}
Que será negativa para $x>1$.

Calculando a primitiva
\begin{align*}
  F & = \int \frac{dx}{x(1+\ln^2(x))}
    \qquad\qquad  u = \ln(x) \qquad du = \frac{dx}{x} \\
    & = \int \frac{du}{1+u^2} \\
    & = \arctan(u) + c \\
    & = \arctan\big(\ln(x)\big) + c
\end{align*}

Calculando a integral imprópria
\begin{align*}
  I
  & = \int_2^\infty \frac{dx}{x(1+\ln^2(x))} \\
  & = \lim_{b\to\infty} \int_2^b \frac{dx}{x(1+\ln^2(x))} \\
  & = \lim_{b\to\infty} \left(\arctan\big(\ln(x)\big)\evalat{2}{b}\right)  \\
  & = \lim_{b\to\infty} \arctan\big(\ln(b)\big) - \arctan\big(\ln(2)\big)  \\
  & = \lim_{t\to\infty} \arctan(t) - \arctan\big(\ln(2)\big)  \\
  & = \frac{\pi}{2} - \arctan\big(\ln(2)\big)  \\
\end{align*}
Como a integral converge a série também converge.
\end{answer}

% 04
%------------------------------------------------------------------------------%
%\vspace{\baselineskip}
\exercise{20}
Sabendo que, para \(n=0, 1, 2, \dots\),
\[
  f^{(n)}(x) = \frac{2(-1)^n n!}{x^{n+1}}
\]
\begin{enumerate}[label=\alph*)]
  \item\ [04] Construa a Série de Taylor de $f(x)$ centrada em 2
  \item\ [04] Determine o intervalo aberto onde a série converge absolutamente
  \item\ [06] Use o polinômio de Taylor de terceiro grau para aproximar \(f(3)\)
  \item\ [06] Estime o erro cometido na aproximação
\end{enumerate}
Não é necessário encontrar a representação decimal dos valores calculados,
mas é necessário realizar as simplificações elementares.

\begin{answer}
\begin{enumerate}[label=\alph*)]
\item
  \[
    f^{(n)}(2)
    = \frac{2(-1)^n n!}{2^{n+1}}
    = \frac{(-1)^n n!}{2^n}
  \]
  \begin{align*}
    T(x)
    & = \sum_{n=0}^{\infty} \frac{f^{(n)}(a)}{n!}(x-1)^n \\
    & = \sum_{n=0}^{\infty} \frac{f^{(n)}(2)}{n!}(x-2)^n \\
    & = \sum_{n=0}^{\infty} \frac{(-1)^n n!}{2^n\,n!}(x-2)^n \\
    & = \sum_{n=0}^{\infty} \frac{(-1)^n}{2^n}(x-2)^n \\
  \end{align*}

\item
  Usando o teste da raiz
  \begin{align*}
    \rho
    & = \lim_{n\to\infty} \sqrt[n]{\abs{a_n}} \\
    & = \lim_{n\to\infty} \sqrt[n]{\frac{\abs{x-2}^n}{2^n}} \\
    & = \lim_{n\to\infty} \frac{\abs{x-2}}{2} \\
    & = \frac{\abs{x-2}}{2} \\
  \end{align*}
  Impondo \(\rho<1\) temos \(\abs{x-2}<2\) portanto \(I=(0, 4)\)

\item
  \[
    P_3(x)
    = \sum_{n=0}^{3} \frac{(-1)^n}{2^n}(x-2)^n
    = 1 - \frac{x-2}{2} + \frac{(x-2)^2}{2^2} - \frac{(x-2)^3}{2^3}
  \]
  \[
    P_3(3)
    = 1 - \frac{3-2}{2} + \frac{(3-2)^2}{2^2} - \frac{(3-2)^3}{2^3}
    = 1 - \frac{1}{2} + \frac{1}{4} - \frac{1}{8}
    = \frac{5}{8}
  \]

\item
  \[
    R_3(x)
    = \frac{f^{(4)}(c)}{4!}(x-2)^4
    = \frac{2(-1)^4 4!}{c^5 4!}(x-2)^4
    = \frac{2(x-2)^4}{c^5}
    \qquad
    c \in [2, x]
  \]
  \[
    R_3(3)
    = \frac{2(3-2)^4}{c^5}
    = \frac{2}{c^5}
    \qquad
    c \in [2, 3]
  \]
  \[
    \abs{R_3(3)}
    \leq \frac{2}{2^5}
    =    \frac{1}{2^4}
    =    \frac{1}{16}
  \]
\end{enumerate}
\end{answer}

%------------------------------------------------------------------------------%
\end{document}
%------------------------------------------------------------------------------%
